\chap{Notation and source concordance}\label{app:concordance}
\subsection{Notation}
\begin{longtable}{@{}p{.25\linewidth}p{.7\linewidth}@{}}
\toprule
Symbol & Meaning\\\midrule\endhead
\(\V,F_v,f^*\) & A declared view category, its fibers, and contravariant restrictions.\\
\(\Gamma(F)\) & The strict set of coherent global families.\\
\(\Gamma_U(F),\Ext_U(s)\) & Partial sections on a full patch and their global extension fiber.\\
\(\delta,\partial,d,D\) & Respectively the view differential, tower differential, a model residual, and a Hilbert defect operator. Each target is declared where used.\\
\(p,p_U,p_L\) & A protected observable, restriction to prescribed views, and a splitting dependent innovation observable.\\
\(\mathfrak o_U,H_U\) & The existence obstruction and the group of invisible coherent changes. Additive and nonabelian conventions are stated separately.\\
\(C_j,J_j,Q_N,S_N\) & Coarsener, chosen lift, restored coarse remainder, and retained finite reconstruction.\\
\(c_f,a_v\) & Coherent nonabelian transition labels and gauge values.\\
\(\mathfrak K_G,J_{\rm comp}\) & The cocycle patch groupoid and its composition-support topology.\\
\(\mathcal E_b(R)\) & Realizations equipped with an actual identification of their boundary with \(b\).\\
\(T_i,K_{ij},R_p\) & Transport to the root, the corresponding root loop operator, and reconstruction from the root.\\
\(\mathcal H,q,P\) & Root residual, an orthogonal kernel projector, and the explicitly defined model repair. A graph repair need not be orthogonal in the global maximum norm.\\
\(\pi_{W,U},H_U\) & The map forgetting part of the prescription and the corresponding invisible coherent ambiguity.\\
\(\phi,\operatorname{EventualHom}\) & A monotone source schedule and transport classes modulo eventual equality.\\
\(\chi_A,M=SA-BC\) & Split exact/residual coordinates and the mismatch between chosen repair lifts.\\
\(\omega,Q_n,C_n,b_n\) & A residual, its prefix-product rate, movement coefficient, and summable step budget. Here $Q_n$ is not the dyadic remainder operator $Q_N$.\\
\(\nu_B,d_W,d_W^{\mathrm{sym}}\) & Residual floor, extended directed transport distance, and its symmetric maximum.\\
\(c(x,y),\kappa(x,y)\) & The two-bit pair phase and its antisymmetric scalar-gauge invariant.\\
\(H,T,r\) & In the budget chapters, $H=\ker p$ with its inherited Hilbert metric, $T=d|_H$, and $r=df$.\\
\(\beta_{T,r}(M),a_\alpha\) & Minimum residual at correction budget $M$, and the quadratic regularized correction.\\
\(C,b,q\) & Closure of the correction image, its projection of the residual, and the irreducible residual floor.\\
\(\simeq,\simeq_L\) & A specified equivalence, and a continuous linear equivalence. Neither notation alone asserts an isometry.\\
\bottomrule
\end{longtable}

\subsection{How to read the correspondence}
The identifiers below are stable editorial identifiers, not counts of new discoveries. A row can combine several source declarations into one exposition. The named declaration, its line range, source hash, and correspondence note are recorded in \code{data/claim_ledger.json} and in the companion Lean guide. Existence of a matching source declaration is not a mechanical proof that every word of an edited statement matches its type. The displayed hypotheses and the printed proof remain part of the review.

\textbf{F} denotes the recovered foundation tree at version 0.20.0. \textbf{M} denotes the focused moving-space tree at version 0.32.0. \textbf{S} denotes the separately recovered Hilbert holonomy source. \textbf{I} denotes a source in the canonical integrated tree, whether inherited or added during the version 1.2.0 formalization. F, M, and S retain provenance, while byte-identical copies of their referenced sources were compiled together in the accepted version 1.2.0 replay of \path{formal/KomorebiSpace/}. All \EditionClaims\ KS catalogue items have source correspondences in that checked closure. The twelve WR extensions are separate written results, listed in Appendix~\ref{app:written}, and have no new Lean counterparts in this edition. This does not assert that every illustrative computation or explanatory sentence is independently formalized.

\input{claim_table}

\par\noindent\begin{minipage}{\linewidth}\subsection{Verification is a separate record}
The foundation archive reports 122 production proof modules, 1,299 named theorem declarations, and 696 indexed constructions. Its recorded strategy includes one complete cumulative source build, followed by a fresh archive replay that reuses a checked 117-module baseline and compiles the five additions and root again. It must not be described as two full clean cumulative builds.

The focused moving-space archive reports 11 production modules and a root, 151 named theorem declarations, and 47 indexed constructions. Its historical replay record reports two complete clean builds of that focused source closure. The archive SHA-256 and the retained source hashes are checked by this publication's integrity script. The focused tree does not replace the whole historical Komorebi Space repository.

The pinned environment in the records is Lean 4.34.0-rc2 with mathlib commit\par\smallskip\noindent\path{aeb72f1e9981b6d35b426b16dca7759be3e7ba06}.\par\smallskip The permitted standard axioms are \code{propext}, \code{Classical.choice}, and \code{Quot.sound}. The compiler, kernel, and cached dependencies remain trusted. For version 1.2.0 the portable runtime was restored and its version and dependency pins checked, but the compiler and mathlib were not rebuilt independently. This expanded publication retains that verified baseline without executing a new compiler replay.

The retained version 1.2.0 formal record covers \EditionFormalModules\ production modules and one root. Clean source replays and an imported axiom audit are separate from archive integrity, exact finite computations, and PDF preflight. The five modules added in that formalization contain \EditionNewTheorems\ explicit named theorems and lemmas, including supporting results. Counts are not claims of mathematical novelty. The complete source hashes, command logs, and baseline scope are in \path{audit/formal/release.json}.

\end{minipage}

\subsection*{Editorial source policy}
Earlier source documents contain historical scope statements and counts. Later checked developments can supersede an earlier statement that a topic is a target. In particular the old outline's stack target is not the status of the later standard stack theorem. This monograph uses the actual relevant declarations rather than treating the earliest outline as a current completion ledger.



\subsection*{Retained source connections and the written extension}
The earlier expansion developed KS59--KS70. KS71--KS80 add exactification, floors, transport costs, and the phase obstruction. Their declarations are listed above. Separately scoped worked models appear in \path{data/deepening_models.json}. WR01--WR12 are indexed separately in \path{data/written_extensions.json}. They have complete printed proofs and exact finite checks, but no new compiled declarations. No new compiler execution is claimed.
